%% bestpaper-writing: figure preamble (see references/figures.md).
%% %% bestpaper-writing: figure preamble (see references/figures.md).
%% %% bestpaper-writing: figure preamble (see references/figures.md).
%% %% bestpaper-writing: figure preamble (see references/figures.md).
%% \input{figures_preamble} after the template's preamble (CVPR kit: at the end of preamble.tex).

\usepackage{tikz}
\usetikzlibrary{positioning,arrows.meta,fit,calc,backgrounds,decorations.pathreplacing,shapes.geometric}

%% Palette: colour-blind safe (Okabe-Ito based). One colour per role, used in every figure.
\definecolor{bpOurs}{RGB}{0,114,178}      % our method / our new component
\definecolor{bpPrior}{RGB}{213,94,0}      % prior work, failure, cost
\definecolor{bpData}{RGB}{0,158,115}      % data, queries, inputs
\definecolor{bpAccent}{RGB}{230,159,0}    % a third role if needed
\definecolor{bpGray}{RGB}{110,110,110}    % reused / frozen / given parts, annotations
\definecolor{bpLight}{RGB}{238,238,238}   % neutral fills
\colorlet{bpOursFill}{bpOurs!18}
\colorlet{bpPriorFill}{bpPrior!15}
\colorlet{bpDataFill}{bpData!15}

%% TikZ styles: boxes by role, one arrow style.
\tikzset{
  bp/base/.style   ={draw=black!70, line width=0.5pt, rounded corners=2pt, align=center,
                     font=\sffamily\scriptsize, minimum height=0.6cm, inner sep=3pt},
  bp/module/.style ={bp/base, fill=white},
  bp/frozen/.style ={bp/base, fill=bpLight, draw=bpGray, text=black!80},
  bp/ours/.style   ={bp/base, fill=bpOursFill, draw=bpOurs, line width=0.7pt},
  bp/prior/.style  ={bp/base, fill=bpPriorFill, draw=bpPrior},
  bp/data/.style   ={bp/base, fill=bpDataFill, draw=bpData},
  bp/todo/.style   ={bp/base, dashed, draw=bpGray, text=bpGray},
  bp/flow/.style   ={-{Latex[length=1.6mm]}, line width=0.6pt, black!75},
  bp/note/.style   ={font=\sffamily\tiny, text=bpGray, align=center},
  bp/panel/.style  ={font=\sffamily\footnotesize\bfseries, anchor=west},
  bp/frame/.style  ={draw=black!70, fill=white, line width=0.4pt},
}

%% A stack of video frames: \bpframes{(x,y)}{width}{height}{number-1}
\newcommand{\bpframes}[4]{%
  \foreach \bpi in {#4,...,0}{\draw[bp/frame] ($#1+(\bpi*0.07,\bpi*0.07)$) rectangle ++(#2,#3);}}

%% Blank area for a figure that needs results; the plan goes into the caption.
\newcommand{\figplaceholder}[2]{%
  \begin{tikzpicture}
    \useasboundingbox (0,0) rectangle (#1,#2);
    \draw[bpGray, dashed, line width=0.6pt, fill=bpLight!45] (0.4pt,0.4pt) rectangle ($(#1,#2)-(0.4pt,0.4pt)$);
    \node[bpGray, font=\sffamily\footnotesize] at ($0.5*(#1,#2)$) {Figure placeholder (see caption)};
  \end{tikzpicture}}

%% The same blank area drawn inside an existing tikzpicture (no nested tikzpicture):
%% \bpplacebox{x}{y}{width}{height} in the picture's units.
\newcommand{\bpplacebox}[4]{%
  \draw[bpGray, dashed, line width=0.6pt, fill=bpLight!45] (#1,#2) rectangle ++(#3,#4);
  \node[bpGray, font=\sffamily\scriptsize, align=center] at ($(#1,#2)+0.5*(#3,#4)$) {Figure placeholder\\(see caption)};}

%% Figure plan inside a caption (red, removed for submission by redefining it as empty).
\providecommand{\TODO}[1]{\textbf{\color{red}[TODO: #1]}}
\newcommand{\FigPlan}[1]{\TODO{Figure plan: #1}}
 after the template's preamble (CVPR kit: at the end of preamble.tex).

\usepackage{tikz}
\usetikzlibrary{positioning,arrows.meta,fit,calc,backgrounds,decorations.pathreplacing,shapes.geometric}

%% Palette: colour-blind safe (Okabe-Ito based). One colour per role, used in every figure.
\definecolor{bpOurs}{RGB}{0,114,178}      % our method / our new component
\definecolor{bpPrior}{RGB}{213,94,0}      % prior work, failure, cost
\definecolor{bpData}{RGB}{0,158,115}      % data, queries, inputs
\definecolor{bpAccent}{RGB}{230,159,0}    % a third role if needed
\definecolor{bpGray}{RGB}{110,110,110}    % reused / frozen / given parts, annotations
\definecolor{bpLight}{RGB}{238,238,238}   % neutral fills
\colorlet{bpOursFill}{bpOurs!18}
\colorlet{bpPriorFill}{bpPrior!15}
\colorlet{bpDataFill}{bpData!15}

%% TikZ styles: boxes by role, one arrow style.
\tikzset{
  bp/base/.style   ={draw=black!70, line width=0.5pt, rounded corners=2pt, align=center,
                     font=\sffamily\scriptsize, minimum height=0.6cm, inner sep=3pt},
  bp/module/.style ={bp/base, fill=white},
  bp/frozen/.style ={bp/base, fill=bpLight, draw=bpGray, text=black!80},
  bp/ours/.style   ={bp/base, fill=bpOursFill, draw=bpOurs, line width=0.7pt},
  bp/prior/.style  ={bp/base, fill=bpPriorFill, draw=bpPrior},
  bp/data/.style   ={bp/base, fill=bpDataFill, draw=bpData},
  bp/todo/.style   ={bp/base, dashed, draw=bpGray, text=bpGray},
  bp/flow/.style   ={-{Latex[length=1.6mm]}, line width=0.6pt, black!75},
  bp/note/.style   ={font=\sffamily\tiny, text=bpGray, align=center},
  bp/panel/.style  ={font=\sffamily\footnotesize\bfseries, anchor=west},
  bp/frame/.style  ={draw=black!70, fill=white, line width=0.4pt},
}

%% A stack of video frames: \bpframes{(x,y)}{width}{height}{number-1}
\newcommand{\bpframes}[4]{%
  \foreach \bpi in {#4,...,0}{\draw[bp/frame] ($#1+(\bpi*0.07,\bpi*0.07)$) rectangle ++(#2,#3);}}

%% Blank area for a figure that needs results; the plan goes into the caption.
\newcommand{\figplaceholder}[2]{%
  \begin{tikzpicture}
    \useasboundingbox (0,0) rectangle (#1,#2);
    \draw[bpGray, dashed, line width=0.6pt, fill=bpLight!45] (0.4pt,0.4pt) rectangle ($(#1,#2)-(0.4pt,0.4pt)$);
    \node[bpGray, font=\sffamily\footnotesize] at ($0.5*(#1,#2)$) {Figure placeholder (see caption)};
  \end{tikzpicture}}

%% The same blank area drawn inside an existing tikzpicture (no nested tikzpicture):
%% \bpplacebox{x}{y}{width}{height} in the picture's units.
\newcommand{\bpplacebox}[4]{%
  \draw[bpGray, dashed, line width=0.6pt, fill=bpLight!45] (#1,#2) rectangle ++(#3,#4);
  \node[bpGray, font=\sffamily\scriptsize, align=center] at ($(#1,#2)+0.5*(#3,#4)$) {Figure placeholder\\(see caption)};}

%% Figure plan inside a caption (red, removed for submission by redefining it as empty).
\providecommand{\TODO}[1]{\textbf{\color{red}[TODO: #1]}}
\newcommand{\FigPlan}[1]{\TODO{Figure plan: #1}}
 after the template's preamble (CVPR kit: at the end of preamble.tex).

\usepackage{tikz}
\usetikzlibrary{positioning,arrows.meta,fit,calc,backgrounds,decorations.pathreplacing,shapes.geometric}

%% Palette: colour-blind safe (Okabe-Ito based). One colour per role, used in every figure.
\definecolor{bpOurs}{RGB}{0,114,178}      % our method / our new component
\definecolor{bpPrior}{RGB}{213,94,0}      % prior work, failure, cost
\definecolor{bpData}{RGB}{0,158,115}      % data, queries, inputs
\definecolor{bpAccent}{RGB}{230,159,0}    % a third role if needed
\definecolor{bpGray}{RGB}{110,110,110}    % reused / frozen / given parts, annotations
\definecolor{bpLight}{RGB}{238,238,238}   % neutral fills
\colorlet{bpOursFill}{bpOurs!18}
\colorlet{bpPriorFill}{bpPrior!15}
\colorlet{bpDataFill}{bpData!15}

%% TikZ styles: boxes by role, one arrow style.
\tikzset{
  bp/base/.style   ={draw=black!70, line width=0.5pt, rounded corners=2pt, align=center,
                     font=\sffamily\scriptsize, minimum height=0.6cm, inner sep=3pt},
  bp/module/.style ={bp/base, fill=white},
  bp/frozen/.style ={bp/base, fill=bpLight, draw=bpGray, text=black!80},
  bp/ours/.style   ={bp/base, fill=bpOursFill, draw=bpOurs, line width=0.7pt},
  bp/prior/.style  ={bp/base, fill=bpPriorFill, draw=bpPrior},
  bp/data/.style   ={bp/base, fill=bpDataFill, draw=bpData},
  bp/todo/.style   ={bp/base, dashed, draw=bpGray, text=bpGray},
  bp/flow/.style   ={-{Latex[length=1.6mm]}, line width=0.6pt, black!75},
  bp/note/.style   ={font=\sffamily\tiny, text=bpGray, align=center},
  bp/panel/.style  ={font=\sffamily\footnotesize\bfseries, anchor=west},
  bp/frame/.style  ={draw=black!70, fill=white, line width=0.4pt},
}

%% A stack of video frames: \bpframes{(x,y)}{width}{height}{number-1}
\newcommand{\bpframes}[4]{%
  \foreach \bpi in {#4,...,0}{\draw[bp/frame] ($#1+(\bpi*0.07,\bpi*0.07)$) rectangle ++(#2,#3);}}

%% Blank area for a figure that needs results; the plan goes into the caption.
\newcommand{\figplaceholder}[2]{%
  \begin{tikzpicture}
    \useasboundingbox (0,0) rectangle (#1,#2);
    \draw[bpGray, dashed, line width=0.6pt, fill=bpLight!45] (0.4pt,0.4pt) rectangle ($(#1,#2)-(0.4pt,0.4pt)$);
    \node[bpGray, font=\sffamily\footnotesize] at ($0.5*(#1,#2)$) {Figure placeholder (see caption)};
  \end{tikzpicture}}

%% The same blank area drawn inside an existing tikzpicture (no nested tikzpicture):
%% \bpplacebox{x}{y}{width}{height} in the picture's units.
\newcommand{\bpplacebox}[4]{%
  \draw[bpGray, dashed, line width=0.6pt, fill=bpLight!45] (#1,#2) rectangle ++(#3,#4);
  \node[bpGray, font=\sffamily\scriptsize, align=center] at ($(#1,#2)+0.5*(#3,#4)$) {Figure placeholder\\(see caption)};}

%% Figure plan inside a caption (red, removed for submission by redefining it as empty).
\providecommand{\TODO}[1]{\textbf{\color{red}[TODO: #1]}}
\newcommand{\FigPlan}[1]{\TODO{Figure plan: #1}}
 after the template's preamble (CVPR kit: at the end of preamble.tex).

\usepackage{tikz}
\usetikzlibrary{positioning,arrows.meta,fit,calc,backgrounds,decorations.pathreplacing,shapes.geometric}

%% Palette: colour-blind safe (Okabe-Ito based). One colour per role, used in every figure.
\definecolor{bpOurs}{RGB}{0,114,178}      % our method / our new component
\definecolor{bpPrior}{RGB}{213,94,0}      % prior work, failure, cost
\definecolor{bpData}{RGB}{0,158,115}      % data, queries, inputs
\definecolor{bpAccent}{RGB}{230,159,0}    % a third role if needed
\definecolor{bpGray}{RGB}{110,110,110}    % reused / frozen / given parts, annotations
\definecolor{bpLight}{RGB}{238,238,238}   % neutral fills
\colorlet{bpOursFill}{bpOurs!18}
\colorlet{bpPriorFill}{bpPrior!15}
\colorlet{bpDataFill}{bpData!15}

%% TikZ styles: boxes by role, one arrow style.
\tikzset{
  bp/base/.style   ={draw=black!70, line width=0.5pt, rounded corners=2pt, align=center,
                     font=\sffamily\scriptsize, minimum height=0.6cm, inner sep=3pt},
  bp/module/.style ={bp/base, fill=white},
  bp/frozen/.style ={bp/base, fill=bpLight, draw=bpGray, text=black!80},
  bp/ours/.style   ={bp/base, fill=bpOursFill, draw=bpOurs, line width=0.7pt},
  bp/prior/.style  ={bp/base, fill=bpPriorFill, draw=bpPrior},
  bp/data/.style   ={bp/base, fill=bpDataFill, draw=bpData},
  bp/todo/.style   ={bp/base, dashed, draw=bpGray, text=bpGray},
  bp/flow/.style   ={-{Latex[length=1.6mm]}, line width=0.6pt, black!75},
  bp/note/.style   ={font=\sffamily\tiny, text=bpGray, align=center},
  bp/panel/.style  ={font=\sffamily\footnotesize\bfseries, anchor=west},
  bp/frame/.style  ={draw=black!70, fill=white, line width=0.4pt},
}

%% A stack of video frames: \bpframes{(x,y)}{width}{height}{number-1}
\newcommand{\bpframes}[4]{%
  \foreach \bpi in {#4,...,0}{\draw[bp/frame] ($#1+(\bpi*0.07,\bpi*0.07)$) rectangle ++(#2,#3);}}

%% Blank area for a figure that needs results; the plan goes into the caption.
\newcommand{\figplaceholder}[2]{%
  \begin{tikzpicture}
    \useasboundingbox (0,0) rectangle (#1,#2);
    \draw[bpGray, dashed, line width=0.6pt, fill=bpLight!45] (0.4pt,0.4pt) rectangle ($(#1,#2)-(0.4pt,0.4pt)$);
    \node[bpGray, font=\sffamily\footnotesize] at ($0.5*(#1,#2)$) {Figure placeholder (see caption)};
  \end{tikzpicture}}

%% The same blank area drawn inside an existing tikzpicture (no nested tikzpicture):
%% \bpplacebox{x}{y}{width}{height} in the picture's units.
\newcommand{\bpplacebox}[4]{%
  \draw[bpGray, dashed, line width=0.6pt, fill=bpLight!45] (#1,#2) rectangle ++(#3,#4);
  \node[bpGray, font=\sffamily\scriptsize, align=center] at ($(#1,#2)+0.5*(#3,#4)$) {Figure placeholder\\(see caption)};}

%% Figure plan inside a caption (red, removed for submission by redefining it as empty).
\providecommand{\TODO}[1]{\textbf{\color{red}[TODO: #1]}}
\newcommand{\FigPlan}[1]{\TODO{Figure plan: #1}}
